\documentclass[11pt,a4paper]{article}
\usepackage[utf8]{inputenc}
\usepackage{amsmath,amssymb}
\usepackage{booktabs}
\usepackage{hyperref}

\title{Normativa de Seguridad Funcional en Robots Colaborativos: ISO 10218 e ISO/TS 15066}
\author{Comité de Seguridad y Robótica Industrial \\
\textit{Centro de Investigación y Transferencia Tecnológica Nova}}
\date{Julio de 2026}

\begin{document}

\maketitle

\begin{abstract}
La coexistencia e interacción cercana entre operadores humanos y manipuladores robóticos sin vallado perimetral físico exige un marco normativo estricto para mitigar riesgos de aprisionamiento, choque o colisión. En este artículo se analiza en profundidad la arquitectura de seguridad requerida bajo las normas internacionales ISO 10218-1/2 (Requisitos de seguridad para robots industriales) y la especificación técnica ISO/TS 15066 (Robots colaborativos). Se detallan los cuatro métodos de operación colaborativa, las fórmulas de distancia de separación protectora en tiempo real y los umbrales biomecánicos de fuerza y presión máxima para contactos transitorios y cuasi-estáticos.
\end{abstract}

\section{Introducción al Concepto Colaborativo}
Un sistema robótico colaborativo (Cobot) no es intrínsecamente seguro por sí mismo; la seguridad funcional depende de la aplicación integral, incluyendo el efector final (gripper), la geometría de la pieza manipulada y el entorno de trabajo compartido. La norma ISO 10218 define el espacio colaborativo como la zona donde robot y operario pueden desempeñar tareas simultáneamente.

\section{Los Cuatro Modos de Operación Colaborativa}
La especificación ISO/TS 15066 estipula cuatro métodos técnicos de operación colaborativa:
\begin{enumerate}
    \item \textbf{Parada Monitoreada con Calificación de Seguridad (SMS):} El robot detiene sus motores manteniendo la energía de los frenos cuando el humano ingresa al espacio de trabajo.
    \item \textbf{Guiado Manual (Hand Guiding):} El operador controla el movimiento del robot mediante una empuñadura ergonómica dotada de un sensor de fuerza-par y un pulsador de validación de 3 posiciones (deadman switch).
    \item \textbf{Monitoreo de Velocidad y Separación (SSM):} Se mantiene una distancia de seguridad dinámica en función de las velocidades relativas detectadas por escáneres láser de seguridad.
    \item \textbf{Limitación de Potencia y Fuerza (PFL):} En caso de contacto inevitable, la transferencia de energía mecánica no supera los umbrales biomecánicos de dolor o lesión.
\end{enumerate}

\section{Monitoreo de Velocidad y Separación: Distancia Protectora $S_p$}
Para garantizar que el robot se detenga por completo antes de que el operario pueda alcanzarlo, la distancia de separación mínima $S_p(t_0)$ en el instante $t_0$ se calcula como:
\begin{equation}
S_p(t_0) = S_h + S_r + S_s + C + Z_d + Z_r
\end{equation}
donde:
\begin{itemize}
    \item $S_h = v_h \cdot (T_r + T_b)$ es la distancia recorrida por el operador con velocidad de aproximación estándar $v_h = 1600\,\text{mm/s}$ (según ISO 13855).
    \item $S_r = v_r \cdot T_r$ es la distancia recorrida por el robot durante el tiempo de reacción del sistema de detección $T_r$.
    \item $S_s$ es la distancia de frenado mecánico propia del robot desde velocidad máxima hasta detención total:
    \begin{equation}
    S_s = \frac{v_r^2}{2 a_{stop}}
    \end{equation}
    \item $C$ es la distancia de penetración intrínseca del escáner láser.
    \item $Z_d$ y $Z_r$ representan tolerancias de medición y posición angular.
\end{itemize}

\section{Límites Biomecánicos de Potencia y Fuerza (PFL)}
La norma ISO/TS 15066 clasifica los contactos físicos en dos categorías:
\begin{itemize}
    \item \textbf{Contacto Cuasi-Estático:} La extremidad del operador queda atrapada entre el efector del robot y una superficie fija rígida.
    \item \textbf{Contacto Transitorio (Dinámico):} El robot golpea una parte del cuerpo libre, disipando energía en el retroceso inercial.
\end{itemize}

\subsection{Tabla de Umbrales Biomecánicos por Región Corporal}
La siguiente tabla resume los valores máximos admisibles de fuerza máxima $F_{max}$ y presión $P_{max}$ antes de inducir dolor o hematoma:

\begin{table}[h]
\centering
\begin{tabular}{lcccc}
\hline
\textbf{Región Corporal} & \textbf{$F_{cuasi}$ [N]} & \textbf{$P_{cuasi}$ [$\text{N/cm}^2$]} & \textbf{$F_{trans}$ [N]} & \textbf{$P_{trans}$ [$\text{N/cm}^2$]} \\
\hline
Cráneo / Frente & 130 & 110 & 260 & 220 \\
Cuello y Garganta & 140 & 140 & 280 & 280 \\
Brazo y Codo & 160 & 180 & 320 & 360 \\
Mano y Dedos & 140 & 190 & 280 & 380 \\
Pecho y Costillas & 140 & 130 & 280 & 260 \\
Muslo y Rodilla & 220 & 220 & 440 & 440 \\
\hline
\end{tabular}
\end{table}

\section{Arquitectura de Control según ISO 13849-1}
Todos los lazos de paro de emergencia, monitoreo de posición segura (SLP), velocidad segura (SLS) y desconexión segura de par (STO) deben cumplir como mínimo con el Nivel de Rendimiento de Seguridad \textbf{PL d, Categoría 3}, garantizando redundancia dual-channel y tolerancia a fallos simples de hardware.

\begin{thebibliography}{9}
\bibitem{iso10218} ISO, \textit{ISO 10218-1:2011 Robots and robotic devices - Safety requirements for industrial robots}, International Organization for Standardization, 2011.
\bibitem{isots15066} ISO, \textit{ISO/TS 15066:2016 Robots and robotic devices - Collaborative robots}, International Organization for Standardization, 2016.
\bibitem{iso13849} ISO, \textit{ISO 13849-1:2015 Safety of machinery - Safety-related parts of control systems}, 2015.
\end{thebibliography}

\end{document}
